\section{Conclusion}
\label{sec:conclusion}

We set out to prove that stable spectral properties indicate model quality --- and proved the exact opposite. The coefficient of variation of participation ratio (CV\_PR), hypothesized to correlate negatively with ImageNet accuracy, instead shows a strong positive correlation ($r = +0.61$, $p < 10^{-10}$). Models with higher accuracy exhibit more variance in randomized SVD estimates, not less.

\subsection{Summary}

In this work, we tested whether SVD estimator variance predicts model quality. Our contributions are:

\begin{enumerate}
    \item \textbf{Methodology validation:} CV\_PR can be reliably extracted from 100+ pretrained models using randomized SVD with 20 seeds (100\% completion, values in $[0.001, 0.033]$).

    \item \textbf{Hypothesis falsification:} The negative correlation hypothesis is decisively rejected. CV\_PR positively correlates with accuracy, with the 95\% confidence interval $[0.506, 0.703]$ excluding all negative values.

    \item \textbf{Interpretive framework:} We identify confounding by model size and richer representations as competing explanations for the unexpected positive correlation.
\end{enumerate}

\subsection{Future Directions}

This falsification opens several research directions grounded in our experimental findings:

\textbf{Confound analysis:} Partial correlation controlling for parameter count would distinguish whether CV\_PR has independent predictive value or merely proxies model complexity.

\textbf{Mechanism investigation:} Why does higher spectral variance accompany better generalization? The inverted causal story --- diversity rather than stability signals quality --- warrants investigation with reformulated hypotheses.

\textbf{Architecture stratification:} Within-family versus cross-family analysis may reveal whether the positive correlation varies by architecture type.

Sometimes the most valuable finding is that our intuitions were wrong. The stability-as-quality assumption, plausible in theory, does not hold for CV\_PR in practice. This negative result redirects research toward understanding why higher variance accompanies better models --- a question that would not have arisen without rigorous hypothesis testing.
